% GENERATED FILE. Edit canonical lessons, then run npm run generate:books.
% canonical-source-hash: fnv1a64:3c2f837b1353f641
% canonical-lessons: ML-C564-alapikkuka, ML-C564-racikkuka, ML-C564-avatarippikkuka, ML-C564-mittuka, ML-C564-abhyasikkuka, ML-R564-review-557-564

\chapter{Composing and Performing}
\label{ch:ml-composing-and-performing}
\hlchaptermodality{564}

\begin{chapteropening}
\textbf{By the end of this chapter:} \emph{I can say to render a song, to compose, to stage, to strum and to train in.}

Complete the last lesson of chapter 564: I can say to render a song, to compose, to stage, to strum and to train in.
\end{chapteropening}

\section[\texorpdfstring{\ml{ആലപിക്കുക} (ālapikkuka)}{ആലപിക്കുക (ālapikkuka)}]{\ml{ആലപിക്കുക} (ālapikkuka) — to render a song; to sing}

\label{lesson:ML-C564-alapikkuka}



\begin{quote}
Before the new one: say the Malayalam for an onlooker, then the Malayalam for applause.
\end{quote}

\subsection*{You'll want to know: \ml{ആലപിക്കുക}}
\textbf{\ml{ആലപിക്കുക}} — \emph{ālapikkuka} — \textquotedblleft{}to render a song; to sing\textquotedblright{}.

\begin{quote}
\ml{ചടങ്ങിൽ} \ml{ഒരു} \ml{കുട്ടി} \ml{പ്രാർത്ഥനാഗാനം} \ml{ആലപിച്ചു}.
\end{quote}

\begin{quote}
\emph{caṭaṅṅil oru kuṭṭi prārtthanāgānaṁ ālapiccu.}
\end{quote}

\begin{quote}
At the ceremony a child sang a prayer song.
\end{quote}

Past: \ml{ആലപിച്ചു}, \emph{ālapiccu}.

Formal, as in programme notes; in speech one says \ml{പാടുക}.

\begin{cousinweb}
\textbf{\ml{ആലപിക്കുക}} (\emph{ālapikkuka}) is built from Sanskrit \emph{ālapati}, \textquotedblleft{}speaks, addresses\textquotedblright{} (CDIAL 1358), with the Malayalam verb ending -\ml{ിക്കുക}. In music \ml{ആലാപനം} is the slow opening of a raga.
\end{cousinweb}

\begin{grammarlens}[title={What you've built}]
One more word for describing what you lived and felt.
\end{grammarlens}

\subsection*{Your turn}
\begin{itemize}
\raggedright
  \item \emph{Say it:} \emph{ālapikkuka}
  \item \emph{Say it:} \emph{ālapikkuka}, once more
  \item \emph{Say it:} \emph{caṭaṅṅil oru kuṭṭi prārtthanāgānaṁ ālapiccu.} — At the ceremony a child sang a prayer song.
  \item \emph{Recall:} say the Malayalam for an onlooker, then the Malayalam for applause, then say \emph{ālapikkuka} again
\end{itemize}

\begin{tcolorbox}[breakable,colback=teal!4,colframe=teal!35!black,title={Before you move on}]
What does \ml{ആലപിക്കുക} mean? (\textquotedblleft{}To render a song; to sing\textquotedblright{}.) Say it once more.
\end{tcolorbox}
\section[\texorpdfstring{\ml{രചിക്കുക} (racikkuka)}{രചിക്കുക (racikkuka)}]{\ml{രചിക്കുക} (racikkuka) — to compose; to author}

\label{lesson:ML-C564-racikkuka}



\begin{quote}
Before the new one: say the Malayalam for applause, then the Malayalam for to render a song.
\end{quote}

\subsection*{You'll want to know: \ml{രചിക്കുക}}
\textbf{\ml{രചിക്കുക}} — \emph{racikkuka} — \textquotedblleft{}to compose; to author\textquotedblright{}.

\begin{quote}
\ml{ആ} \ml{കവി} \ml{ഒരുപാട്} \ml{ഗാനങ്ങൾ} \ml{രചിച്ചു}.
\end{quote}

\begin{quote}
\emph{ā kavi orupāṭŭ gānaṅṅaḷ raciccu.}
\end{quote}

\begin{quote}
That poet composed many songs.
\end{quote}

Past: \ml{രചിച്ചു}, \emph{raciccu}.

The noun is \ml{രചന}, a composition; its author is \ml{രചയിതാവ്}.

\begin{cousinweb}
\textbf{\ml{രചിക്കുക}} (\emph{racikkuka}) is built from Sanskrit \emph{racana}, \textquotedblleft{}making, arranging, composing\textquotedblright{} (Monier-Williams), with the Malayalam verb ending -\ml{ിക്കുക}.
\end{cousinweb}

\begin{grammarlens}[title={What you've built}]
One more word for describing what you lived and felt.
\end{grammarlens}

\subsection*{Your turn}
\begin{itemize}
\raggedright
  \item \emph{Say it:} \emph{racikkuka}
  \item \emph{Say it:} \emph{racikkuka}, once more
  \item \emph{Say it:} \emph{ā kavi orupāṭŭ gānaṅṅaḷ raciccu.} — That poet composed many songs.
  \item \emph{Recall:} say the Malayalam for applause, then the Malayalam for to render a song, then say \emph{racikkuka} again
\end{itemize}

\begin{tcolorbox}[breakable,colback=teal!4,colframe=teal!35!black,title={Before you move on}]
What does \ml{രചിക്കുക} mean? (\textquotedblleft{}To compose; to author\textquotedblright{}.) Say it once more.
\end{tcolorbox}
\section[\texorpdfstring{\ml{അവതരിപ്പിക്കുക} (avatarippikkuka)}{അവതരിപ്പിക്കുക (avatarippikkuka)}]{\ml{അവതരിപ്പിക്കുക} (avatarippikkuka) — to stage; to present}

\label{lesson:ML-C564-avatarippikkuka}



\begin{quote}
Before the new one: say the Malayalam for to render a song, then the Malayalam for to compose.
\end{quote}

\subsection*{You'll want to know: \ml{അവതരിപ്പിക്കുക}}
\textbf{\ml{അവതരിപ്പിക്കുക}} — \emph{avatarippikkuka} — \textquotedblleft{}to stage; to present\textquotedblright{}.

\begin{quote}
\ml{കുട്ടികൾ} \ml{വേദിയിൽ} \ml{ഒരു} \ml{നാടകം} \ml{അവതരിപ്പിച്ചു}.
\end{quote}

\begin{quote}
\emph{kuṭṭikaḷ vēdiyil oru nāṭakaṁ avatarippiccu.}
\end{quote}

\begin{quote}
The children put on a play on the stage.
\end{quote}

Past: \ml{അവതരിപ്പിച്ചു}, \emph{avatarippiccu}.

Also to present a report or a programme; a presenter is \ml{അവതാരകൻ}.

\begin{cousinweb}
\textbf{\ml{അവതരിപ്പിക്കുക}} (\emph{avatarippikkuka}) is built on Sanskrit \emph{avatāra}, \textquotedblleft{}coming down, descent\textquotedblright{} (CDIAL 777), with the Malayalam causative ending -\ml{ിപ്പിക്കുക}.
\end{cousinweb}

\begin{grammarlens}[title={What you've built}]
One more word for describing what you lived and felt.
\end{grammarlens}

\subsection*{Your turn}
\begin{itemize}
\raggedright
  \item \emph{Say it:} \emph{avatarippikkuka}
  \item \emph{Say it:} \emph{avatarippikkuka}, once more
  \item \emph{Say it:} \emph{kuṭṭikaḷ vēdiyil oru nāṭakaṁ avatarippiccu.} — The children put on a play on the stage.
  \item \emph{Recall:} say the Malayalam for to render a song, then the Malayalam for to compose, then say \emph{avatarippikkuka} again
\end{itemize}

\begin{tcolorbox}[breakable,colback=teal!4,colframe=teal!35!black,title={Before you move on}]
What does \ml{അവതരിപ്പിക്കുക} mean? (\textquotedblleft{}To stage; to present\textquotedblright{}.) Say it once more.
\end{tcolorbox}
\section[\texorpdfstring{\ml{മീട്ടുക} (mīṭṭuka)}{മീട്ടുക (mīṭṭuka)}]{\ml{മീട്ടുക} (mīṭṭuka) — to strum; to pluck (strings)}

\label{lesson:ML-C564-mittuka}



\begin{quote}
Before the new one: say the Malayalam for to compose, then the Malayalam for to stage.
\end{quote}

\subsection*{You'll want to know: \ml{മീട്ടുക}}
\textbf{\ml{മീട്ടുക}} — \emph{mīṭṭuka} — \textquotedblleft{}to strum; to pluck (strings)\textquotedblright{}.

\begin{quote}
\ml{അവൻ} \ml{ഗിറ്റാറിന്റെ} \ml{കമ്പികൾ} \ml{പതുക്കെ} \ml{മീട്ടി}.
\end{quote}

\begin{quote}
\emph{avan giṟṟāṟinṟe kampikaḷ patukke mīṭṭi.}
\end{quote}

\begin{quote}
He gently strummed the strings of the guitar.
\end{quote}

Past: \ml{മീട്ടി}, \emph{mīṭṭi}.

Said of strings; a drum is \ml{കൊട്ടുക}.

\begin{cousinweb}
\textbf{\ml{മീട്ടുക}} (\emph{mīṭṭuka}) is native Dravidian: DEDR 4848 gives it as to tap or strike, as the strings of an instrument (DEDR 4848).
\end{cousinweb}

\begin{grammarlens}[title={What you've built}]
One more word for describing what you lived and felt.
\end{grammarlens}

\subsection*{Your turn}
\begin{itemize}
\raggedright
  \item \emph{Say it:} \emph{mīṭṭuka}
  \item \emph{Say it:} \emph{mīṭṭuka}, once more
  \item \emph{Say it:} \emph{avan giṟṟāṟinṟe kampikaḷ patukke mīṭṭi.} — He gently strummed the strings of the guitar.
  \item \emph{Recall:} say the Malayalam for to compose, then the Malayalam for to stage, then say \emph{mīṭṭuka} again
\end{itemize}

\begin{tcolorbox}[breakable,colback=teal!4,colframe=teal!35!black,title={Before you move on}]
What does \ml{മീട്ടുക} mean? (\textquotedblleft{}To strum; to pluck (strings)\textquotedblright{}.) Say it once more.
\end{tcolorbox}
\section[\texorpdfstring{\ml{അഭ്യസിക്കുക} (abhyasikkuka)}{അഭ്യസിക്കുക (abhyasikkuka)}]{\ml{അഭ്യസിക്കുക} (abhyasikkuka) — to train in; to learn (an art)}

\label{lesson:ML-C564-abhyasikkuka}



\begin{quote}
Before the new one: say the Malayalam for to stage, then the Malayalam for to strum.
\end{quote}

\subsection*{You'll want to know: \ml{അഭ്യസിക്കുക}}
\textbf{\ml{അഭ്യസിക്കുക}} — \emph{abhyasikkuka} — \textquotedblleft{}to train in; to learn (an art)\textquotedblright{}.

\begin{quote}
\ml{അവൾ} \ml{പത്ത്} \ml{വർഷം} \ml{നൃത്തം} \ml{അഭ്യസിച്ചു}.
\end{quote}

\begin{quote}
\emph{avaḷ pattŭ varṣaṁ nṛttaṁ abhyasiccu.}
\end{quote}

\begin{quote}
She trained in dance for ten years.
\end{quote}

Past: \ml{അഭ്യസിച്ചു}, \emph{abhyasiccu}.

The noun is \ml{അഭ്യാസം}, practice, also a trick or a stunt.

\begin{cousinweb}
\textbf{\ml{അഭ്യസിക്കുക}} (\emph{abhyasikkuka}) is built from Sanskrit \emph{abhyāsa}, \textquotedblleft{}repetition, practice\textquotedblright{} (Monier-Williams), with the Malayalam verb ending -\ml{ിക്കുക}.
\end{cousinweb}

\begin{grammarlens}[title={What you've built}]
One more word for describing what you lived and felt.
\end{grammarlens}

\subsection*{Your turn}
\begin{itemize}
\raggedright
  \item \emph{Say it:} \emph{abhyasikkuka}
  \item \emph{Say it:} \emph{abhyasikkuka}, once more
  \item \emph{Say it:} \emph{avaḷ pattŭ varṣaṁ nṛttaṁ abhyasiccu.} — She trained in dance for ten years.
  \item \emph{Recall:} say the Malayalam for to stage, then the Malayalam for to strum, then say \emph{abhyasikkuka} again
\end{itemize}

\begin{tcolorbox}[breakable,colback=teal!4,colframe=teal!35!black,title={Before you move on}]
What does \ml{അഭ്യസിക്കുക} mean? (\textquotedblleft{}To train in; to learn (an art)\textquotedblright{}.) Say it once more.
\end{tcolorbox}
\section[(dialogue)]{Words from chapters 557–564 — a review}

\label{lesson:ML-R564-review-557-564}



\begin{quote}
Say the Malayalam for to hand-feed and to scold.
\end{quote}

\subsection*{Your turn}
Say these aloud:

\begin{itemize}
\raggedright
  \item to hand-feed, to scold, to pamper, to cuddle, a get-together
  \item pleasantries, early childhood, to prattle, grown-up, a twin
  \item to look after, to hold close, to get along, to cherish, to ask forgiveness
  \item an excursion, painting, a sculpture, a sculptor, an exhibition
  \item a concert, a musical instrument, to beat a drum, a guitar, a flute
  \item rhythm, a raga, a director, a screenplay, a novel
  \item a story character, a stage scene, a stage, an onlooker, applause
  \item to render a song, to compose, to stage, to strum, to train in
\end{itemize}

\begin{tcolorbox}[breakable,colback=teal!4,colframe=teal!35!black,title={Before you move on}]
To hand-feed? (\textbf{\ml{ഊട്ടുക}}, \emph{ūṭṭuka}.) To scold? (\textbf{\ml{ശാസിക്കുക}}, \emph{śāsikkuka}.) To pamper? (\textbf{\ml{ലാളിക്കുക}}, \emph{lāḷikkuka}.) To cuddle? (\textbf{\ml{താലോലിക്കുക}}, \emph{tālōlikkuka}.) A get-together? (\textbf{\ml{ഒത്തുചേരൽ}}, \emph{ottucēral}.) Pleasantries? (\textbf{\ml{കുശലം}}, \emph{kuśalaṁ}.) Early childhood? (\textbf{\ml{ബാല്യം}}, \emph{bālyaṁ}.) To prattle? (\textbf{\ml{കൊഞ്ചുക}}, \emph{koñcuka}.) Grown-up? (\textbf{\ml{മുതിർന്ന}}, \emph{mutirnna}.) A twin? (\textbf{\ml{ഇരട്ടക്കുട്ടി}}, \emph{iraṭṭakkuṭṭi}.) To look after? (\textbf{\ml{പരിപാലിക്കുക}}, \emph{paripālikkuka}.) To hold close? (\textbf{\ml{ചേർത്തുപിടിക്കുക}}, \emph{cērttupiṭikkuka}.) To get along? (\textbf{\ml{ഒത്തുപോകുക}}, \emph{ottupōkuka}.) To cherish? (\textbf{\ml{കാത്തുസൂക്ഷിക്കുക}}, \emph{kāttusūkṣikkuka}.) To ask forgiveness? (\textbf{\ml{ക്ഷമ} \ml{ചോദിക്കുക}}, \emph{kṣama cōdikkuka}.) An excursion? (\textbf{\ml{വിനോദയാത്ര}}, \emph{vinōdayātra}.) Painting? (\textbf{\ml{ചിത്രരചന}}, \emph{citraracana}.) A sculpture? (\textbf{\ml{ശില്പം}}, \emph{śilpaṁ}.) A sculptor? (\textbf{\ml{ശില്പി}}, \emph{śilpi}.) An exhibition? (\textbf{\ml{പ്രദർശനം}}, \emph{pradarśanaṁ}.) A concert? (\textbf{\ml{കച്ചേരി}}, \emph{kaccēri}.) A musical instrument? (\textbf{\ml{വാദ്യം}}, \emph{vādyaṁ}.) To beat a drum? (\textbf{\ml{കൊട്ടുക}}, \emph{koṭṭuka}.) A guitar? (\textbf{\ml{ഗിറ്റാർ}}, \emph{giṟṟār}.) A flute? (\textbf{\ml{ഓടക്കുഴൽ}}, \emph{ōṭakkuḻal}.) Rhythm? (\textbf{\ml{താളം}}, \emph{tāḷaṁ}.) A raga? (\textbf{\ml{രാഗം}}, \emph{rāgaṁ}.) A director? (\textbf{\ml{സംവിധായകൻ}}, \emph{saṁvidhāyakan}.) A screenplay? (\textbf{\ml{തിരക്കഥ}}, \emph{tirakkatha}.) A novel? (\textbf{\ml{നോവൽ}}, \emph{nōval}.) A story character? (\textbf{\ml{കഥാപാത്രം}}, \emph{kathāpātraṁ}.) A stage scene? (\textbf{\ml{രംഗം}}, \emph{raṁgaṁ}.) A stage? (\textbf{\ml{വേദി}}, \emph{vēdi}.) An onlooker? (\textbf{\ml{കാണി}}, \emph{kāṇi}.) Applause? (\textbf{\ml{കയ്യടി}}, \emph{kayyaṭi}.) To render a song? (\textbf{\ml{ആലപിക്കുക}}, \emph{ālapikkuka}.) To compose? (\textbf{\ml{രചിക്കുക}}, \emph{racikkuka}.) To stage? (\textbf{\ml{അവതരിപ്പിക്കുക}}, \emph{avatarippikkuka}.) To strum? (\textbf{\ml{മീട്ടുക}}, \emph{mīṭṭuka}.) To train in? (\textbf{\ml{അഭ്യസിക്കുക}}, \emph{abhyasikkuka}.)
\end{tcolorbox}
